\documentclass[10pt,reqno]{amsart}

\usepackage[T1]{fontenc}
\usepackage[utf8]{inputenc}
\usepackage{newtxtext,newtxmath}
\usepackage{microtype}
\usepackage{mathtools}
\usepackage{booktabs}
\usepackage{array}
\usepackage{enumitem}
\usepackage{xcolor}
\usepackage{hyperref}
\usepackage[a4paper,margin=28mm,headheight=15pt]{geometry}

\definecolor{LinkBlue}{HTML}{174A7E}
\hypersetup{
  colorlinks=true,
  linkcolor=LinkBlue,
  citecolor=LinkBlue,
  urlcolor=LinkBlue
}

\setlength{\parskip}{2pt}
\setlist[itemize]{leftmargin=6mm,itemsep=1.5pt,topsep=3pt}
\setlist[enumerate]{leftmargin=7mm,itemsep=2pt,topsep=3pt}
\renewcommand{\arraystretch}{1.15}
\emergencystretch=2em

\newtheorem{theorem}{Theorem}[section]
\newtheorem{proposition}[theorem]{Proposition}
\newtheorem{lemma}[theorem]{Lemma}
\newtheorem{corollary}[theorem]{Corollary}
\theoremstyle{definition}
\newtheorem{definition}[theorem]{Definition}
\newtheorem{remark}[theorem]{Remark}

\title[Non-Stabilization of Chow Semistable Pluricanonical Limits]
{Non-Stabilization of Chow Semistable Pluricanonical Limits for Smooth Hypersurfaces}
\author{Argus Open Research Project}
\date{August 27, 2026}
\subjclass[2020]{Primary 14J10; Secondary 14L24, 32Q20}
\begin{document}
\begin{abstract}
We record the theorem of Wang and Xu (2014) disproving the stabilization of Chow semistable $r$-canonical limits as $r\to\infty$ for degenerating one-parameter families of canonically polarized varieties, including smooth hypersurfaces of large degree in $\mathbb{P}^3$.
\end{abstract}
\maketitle
\section{Main Refutation}
\begin{theorem}[Wang--Xu, 2014]
There exist flat projective families of canonically polarized manifolds over a punctured curve (including degenerations of smooth surfaces of degree $d\gg 1$ in $\mathbb{P}^3$) whose asymptotic Chow semistable limits with respect to $r K_X$ do not stabilize as $r\to\infty$ whenever the K-semistable limit differs from the asymptotic Chow limit.
\end{theorem}
\end{document}
